\documentclass[12pt]{article}
\newcommand{\DocShort}{Incremento incremental}
\input{preamble_population.tex}
\begin{document}
\doccover{Método del incremento incremental}{2026-05}
\handoutintro{Proyectar la población considerando tanto el incremento medio como el incremento incremental medio.}{Ciudades de tamaño medio cuyo crecimiento viene acelerándose de manera sostenida, pero no explosiva.}

\section{Datos ejemplo}
Se usará una serie con períodos censales de $10$ años. En este ejemplo, tanto el incremento medio $X$ como el incremento incremental medio $Y$ se expresarán por período censal.

\begin{center}
\begin{tabular}{lcccccc}
\toprule
Año & 1970 & 1980 & 1990 & 2000 & 2010 & 2020 \\
\midrule
Población & 68,400 & 78,200 & 89,600 & 103,900 & 120,700 & 140,400 \\
\bottomrule
\end{tabular}
\end{center}

\section{Ecuaciones mínimas}
El incremento medio por período censal es:
\begin{equation}
X=\frac{1}{N}\sum_{i=1}^{N}\Delta P_i
\end{equation}
donde $\Delta P_i$ es el incremento de población del período censal $i$ y $N$ es el número de períodos.

El incremento incremental medio por período censal es:
\begin{equation}
Y=\frac{1}{N-1}\sum_{i=2}^{N}\left(\Delta P_i-\Delta P_{i-1}\right)
\end{equation}
donde $Y$ mide cuánto aumenta, en promedio, el incremento de un período al siguiente.

La población proyectada después de $n$ períodos es:
\begin{equation}
P_n=P_{\text{actual}}+nX+\frac{n(n+1)}{2}Y
\end{equation}
donde $P_{\text{actual}}$ es la población del último censo.

\weightednote{Si los intervalos censales no son regulares, primero conviene convertir los incrementos a una base común, por ejemplo habitantes por año, y luego obtener los promedios ponderados con pesos iguales a la duración de cada período en años.}

\section{Procedimiento}
\begin{enumerate}
\item Calcular los incrementos por período censal.
\item Calcular los incrementos de los incrementos.
\item Obtener el promedio de ambos conjuntos.
\item Sustituir en la ecuación para cada período futuro.
\end{enumerate}

\section{Ejemplo resuelto}
\begin{center}
\begin{tabular}{crrr}
\toprule
Año & Población & Incremento del período & Incremento incremental \\
\midrule
1970 & 68,400 & -- & -- \\
1980 & 78,200 & 9,800 & -- \\
1990 & 89,600 & 11,400 & 1,600 \\
2000 & 103,900 & 14,300 & 2,900 \\
2010 & 120,700 & 16,800 & 2,500 \\
2020 & 140,400 & 19,700 & 2,900 \\
\midrule
Promedio & & 14,400 & 2,475 \\
\bottomrule
\end{tabular}
\end{center}

\[
P_{2030}=\num{140400}+1(\num{14400})+\frac{1(1+1)}{2}(\num{2475})=\num{157275}
\]
\[
P_{2040}=\num{140400}+2(\num{14400})+\frac{2(2+1)}{2}(\num{2475})=\num{176625}
\]
\[
P_{2050}=\num{140400}+3(\num{14400})+\frac{3(3+1)}{2}(\num{2475})=\num{198450}
\]

\resultbox{La población proyectada es 157,275 habitantes en 2030, 176,625 en 2040 y 198,450 en 2050.}

\section{Teoría breve}
Es una corrección del método aritmético. Mantiene una parte lineal, pero agrega un término que incorpora la aceleración media del crecimiento.

\section{Observaciones de uso}
\begin{itemize}
\item Suele producir valores intermedios entre el método aritmético y el geométrico.
\item Conviene cuando la ciudad ha mostrado incrementos crecientes.
\item Debe aclararse si $X$ y $Y$ están expresados por año, por década o por otro período.
\end{itemize}
\end{document}
